% ============================================================
% Week 2 — 09 Mar – 13 Mar 2026
% Compile:  cd report && make week_02
% ============================================================
\documentclass[11pt, a4paper]{article}

\usepackage[margin=2.5cm]{geometry}
\usepackage{graphicx}
\usepackage{booktabs}
\usepackage{multirow}
\usepackage{array}
\usepackage{xcolor}
\usepackage{hyperref}
\usepackage{amsmath}
\usepackage{amssymb}
\usepackage{caption}
\usepackage{subcaption}
\usepackage{float}
\usepackage{enumitem}
\usepackage{fancyhdr}
\usepackage{titlesec}
\usepackage{parskip}

\definecolor{sota}{HTML}{2E7D32}
\definecolor{good}{HTML}{1565C0}
\definecolor{warn}{HTML}{E65100}

\setlength{\headheight}{14pt}
\pagestyle{fancy}
\fancyhf{}
\lhead{\textbf{ESS Research} --- Week 2 Update}
\rhead{Annotation-Free Surgical Instrument Segmentation}
\cfoot{\thepage}

\titleformat{\section}{\large\bfseries}{}{0em}{}[\titlerule]
\titleformat{\subsection}{\normalsize\bfseries}{}{0em}{}

\hypersetup{colorlinks=true, linkcolor=blue, urlcolor=blue}
\newcolumntype{C}[1]{>{\centering\arraybackslash}p{#1}}

\begin{document}

% ── Cover ────────────────────────────────────────────────────
\begin{center}
    {\LARGE \textbf{Towards Annotation-Free Surgical Instrument Segmentation}}\\[0.3em]
    {\large via Foundation Model Adaptation and Robot Kinematics}\\[0.5em]
    \rule{0.5\linewidth}{0.4pt}\\[0.4em]
    {\normalsize Weekly Research Progress Report --- \textbf{Week 2}}\\[0.3em]
    {\normalsize Week dates: 09 Mar -- 13 Mar 2026 \quad|\quad Report date: 13 March 2026}\\[0.3em]
    {\normalsize \href{https://github.com/skhan61/ess-research}{\texttt{github.com/skhan61/ess-research}}}
\end{center}
\vspace{0.5em}

\noindent Week~1 established that SAM3+LoRA+box exceeds task-specific SOTA and that
spatial prompt quality is the primary performance bottleneck.
Week~2 records the research thinking that follows: \textbf{what does the
bottleneck reveal about the deeper problem, and where does it point within the
broader space of robotic perception?}
No new experiments were run. This report documents an open research direction.

\hrule
\vspace{1em}

% ============================================================
\section{What Week 1 Actually Showed}
% ============================================================

The prompt ablation (H4) produced a result that deserves closer reading.
The box prompt --- derived directly from the GT mask --- is the strongest
prompt. This creates a circular dependency: the best spatial prior comes
from the annotation we are trying to replace.
The practical implication is that the performance ceiling observed in Week~1
is set not by the model but by the annotation quality itself.
If the annotations contain error, the reported Dice contains that error,
and the model has no way to exceed it.

This raises a more fundamental question than ``which prompt is best'':
\textbf{can we trust the GT masks we are training and evaluating against?}

% ============================================================
\section{The Ground Truth Reliability Problem}
% ============================================================

Manual pixel-level annotation of surgical instruments introduces error that
is both irreducible and unmeasurable. Annotator disagreement on fine instrument
boundaries is typically 3--10 pixels. There is no reference standard against
which to quantify this error. As a result, Dice scores computed against manual
GT describe similarity to a human drawing, not accuracy of true instrument
boundary delineation. This is a known but largely unaddressed problem in
surgical vision benchmarking.

The question becomes: is there a GT source whose error is \textit{bounded}
and \textit{physically interpretable}?

% ============================================================
\section{Robot Kinematics as a GT Source}
% ============================================================

\subsection{The Core Idea}

In a robotic surgical system, the instrument pose --- base and tip positions
in 3D --- is known from forward kinematics and encoder readings.
Projecting this onto the image plane yields a geometric instrument mask
whose boundary error is bounded by the encoder precision, not by human
judgement.

\subsection{How GT is Generated}

The pipeline from robot data to GT mask involves three steps:

\begin{enumerate}[noitemsep]
    \item \textbf{Forward kinematics:} robot joint encoder readings
          $\to$ 3D positions $P_\text{base},\, P_\text{tip}$ in robot
          base frame.
    \item \textbf{Hand-eye projection:} apply the robot-to-camera
          transform $T_\text{cam}^{\text{robot}}$ (hand-eye calibration)
          $\to$ 2D image coordinates
          $p_\text{base} = \pi(T \cdot P_\text{base})$,\;
          $p_\text{tip}  = \pi(T \cdot P_\text{tip})$,
          where $\pi(\cdot)$ is the camera projection.
    \item \textbf{Mask generation:} for each image pixel $x$, assign
          instrument label if its distance to the projected centerline
          segment is less than the known instrument radius $r$:
\end{enumerate}

\begin{equation}
    M^*(x) = \mathbf{1}\!\left[
        \min_{s\in[0,1]} \left\| x -
        \bigl[(1{-}s)\,p_{\text{base}} + s\,p_{\text{tip}}\bigr]
        \right\|_2 < r
    \right]
\end{equation}

The error in $M^*$ propagates from two sources: encoder precision
(Step~1) and hand-eye calibration accuracy (Step~2). Both are
measurable. The lab dataset carries $\leq 1$\,mm combined accuracy,\footnote{Based
on email correspondence and meeting discussion with Prof.\ Yangming Lee
(Week~1--2). Explicit confirmation of data format and available metadata
is pending.} replacing an \textit{unknown} annotation error with a
\textit{known} calibration error.

\subsection{What This Means for Evaluation}

A Dice score computed against $M^*$ has a direct physical interpretation:
it measures how close the predicted mask is to the true instrument boundary,
within the tolerance of the robot's own measurement accuracy.
This is qualitatively different from Dice against manual GT.

\subsection{The Deployment Implication}

Kinematic data is available in a robotic system during data collection.
It is \textit{not} available in general clinical deployment: non-robotic
procedures, different hospitals, different instruments.
The natural architecture is therefore:

\begin{itemize}[noitemsep]
    \item \textbf{Training:} use kinematic data to generate $M^*$;
          fine-tune SAM3+LoRA on these masks.
    \item \textbf{Inference:} image + text prompt only. No robot.
          No sensor. No annotation.
\end{itemize}

The model acquires its knowledge of instrument appearance from robot-supervised
training and transfers it to settings where the robot is absent.

% ============================================================
\section{Data Requirements for the Full Pipeline}
% ============================================================

\subsection{What Is Needed}

\begin{center}
\begin{tabular}{p{4cm} p{4.5cm} p{4cm}}
\toprule
\textbf{Data} & \textbf{Used For} & \textbf{When Needed} \\
\midrule
Video frames (RGB + timestamps) &
    Model input &
    Training + inference \\
Joint encoder readings \textit{or} pre-computed $(P_\text{base}, P_\text{tip})$ &
    Forward kinematics $\to$ 3D pose &
    Training only \\
Hand-eye calibration matrix $T_\text{cam}^\text{robot}$ &
    Project 3D pose $\to$ image coordinates &
    Training only \\
Camera intrinsics $K$ + distortion &
    Camera projection $\pi(\cdot)$ &
    Training only \\
Instrument radius $r$ (per instrument type) &
    Mask generation via Eq.~(1) &
    Training only \\
Frame--pose timestamp sync &
    Align image to correct pose &
    Training only \\
\midrule
\textbf{At inference} & Image + text prompt only & \textbf{Nothing else} \\
\bottomrule
\end{tabular}
\end{center}

The $\leq 1$\,mm accuracy stated for the lab dataset implies that
$T_\text{cam}^\text{robot}$ and $K$ were already computed and validated
as part of data collection. These are therefore available from the dataset
owner alongside the kinematic and video data.

% ============================================================
\section{Where This Sits in Robotic Perception}
% ============================================================

The broader question this belongs to is:
\textbf{how does a robot use its own proprioceptive knowledge to supervise
its visual perception?}

In manipulation and navigation, this is an active research area --- using
internal state (joint angles, odometry, force readings) to label visual
data without human annotation. Surgical robotics is a natural domain for
this because the robot already records precise kinematic data as part of
normal operation.

The connection to segmentation is one instance of this.
A more general formulation: the robot's kinematic state at time $t$ provides
a spatial prior over the image at time $t$. Any visual learning task that
benefits from knowing where the robot's end-effector is --- segmentation,
depth estimation, tissue contact detection, scene reconstruction ---
could in principle be supervised this way.

Whether segmentation specifically is the right task to anchor a paper,
or whether the contribution is better framed at the level of the
kinematic-visual supervision principle, remains an open question.

% ============================================================
\section{What Prior Work Has and Has Not Done}
% ============================================================

\begin{center}
\begin{tabular}{p{3cm} C{1.2cm} p{4.5cm} p{4cm}}
\toprule
\textbf{Work} & \textbf{Year} & \textbf{Dataset} & \textbf{Limitation} \\
\midrule
Rocha \& Padoy~\cite{rocha2019kinematic} & 2019 &
Custom dVRK. Laparoscopic frames + kinematics. Not public. &
Kinematics needed at inference. No foundation model. \\
CaRTS~\cite{ding2022carts} & 2022 &
dVRK synthetic + real; smoke, blood, altered backgrounds. &
Kinematics needed at inference. Robot required at deployment. \\
SAF-IS~\cite{sestini2025safis} & 2025 &
EndoVis 2017 \& 2018 (public). Annotation-free via unsupervised masks. &
No kinematic GT. No physically bounded error. \\
\bottomrule
\end{tabular}
\end{center}

\noindent Prior work that uses kinematic GT \cite{rocha2019kinematic,ding2022carts}
requires the robot at inference because pre-SAM models were too weak to
operate on visual input alone. SAM3 closes this gap.
Prior work that achieves annotation-free inference \cite{sestini2025safis} does not use
kinematics and therefore cannot provide a physically bounded GT error.
The combination --- kinematic GT at training, foundation model inference ---
has not been done.

% ============================================================
\section{Next Steps}
% ============================================================

The lab dataset contains video frames and robot pose data.
\textbf{There are no existing GT masks.} The $\leq 1$\,mm figure refers
to the robot encoder accuracy --- meaning we can derive GT masks from
the pose at that accuracy. This derivation is the contribution.

Assuming dataset access:

\begin{enumerate}[noitemsep]
    \item Apply Eq.~(1) to robot pose data to generate $M^*$ for all
          frames --- creating GT masks that do not otherwise exist.
    \item Fine-tune the existing SAM3+LoRA checkpoint on these
          kinematic-derived masks.
    \item Evaluate on held-out frames: does the model trained on
          derived GT segment correctly?
    \item Characterise the result: segmentation accuracy is now
          bounded by $\leq 1$\,mm encoder precision rather than
          by annotation quality.
\end{enumerate}

% ============================================================
\section{Repository}
% ============================================================

\begin{center}
\begin{tabular}{ll}
\toprule
\textbf{Item} & \textbf{Detail} \\
\midrule
Repository   & \texttt{github.com/skhan61/ess-research} \\
Entry point  & \texttt{uv run python scripts/run\_all\_experiments.py --plan research\_plan.yaml} \\
Debug flag   & \texttt{--debug} \\
\bottomrule
\end{tabular}
\end{center}

\nocite{qin2020uw,hfrfnet2025,ravi2024sam2,yue2024surgicalsam,wu2024pointtracking,liu2024surgsam2}
\bibliographystyle{plain}
\bibliography{../references}

\end{document}
